\documentclass[11pt,oneside]{report}
\usepackage[T1]{fontenc}
\usepackage[utf8]{inputenc}
\usepackage{lmodern,geometry,microtype,xcolor,hyperref,enumitem,booktabs,longtable,array,fancyhdr,listings}
\geometry{margin=0.9in}
\definecolor{napariblue}{HTML}{2463A6}
\hypersetup{colorlinks=true,linkcolor=napariblue,urlcolor=napariblue,pdftitle={napari-raman-widget Detailed Guide}}
\setlength{\parindent}{0pt}
\setlength{\parskip}{0.5em}
\setlist{nosep,leftmargin=*}
\pagestyle{fancy}\fancyhf{}\fancyhead[L]{\textit{napari-raman-widget}}\fancyhead[R]{\thepage}
\newcommand{\ui}[1]{\textbf{#1}}
\newcommand{\code}[1]{\texttt{#1}}
\newcommand{\fieldtable}{\begin{longtable}{@{}>{\raggedright\arraybackslash}p{0.25\textwidth}>{\raggedright\arraybackslash}p{0.18\textwidth}>{\raggedright\arraybackslash}p{0.50\textwidth}@{}}\toprule\textbf{Control} & \textbf{Default / allowed} & \textbf{Meaning and use}\\\midrule\endhead}
\newcommand{\endfieldtable}{\bottomrule\end{longtable}}

\title{\textbf{napari-raman-widget}\\\Large Detailed User and Architecture Guide}
\author{Raman Microscopy Software Documentation}
\date{\today}

\begin{document}
\maketitle
\thispagestyle{empty}
\begin{abstract}
This guide documents the current controls in the napari Raman sidebar. It explains what each editable field does,
which napari layer is consumed, what must be prepared first, and which output is
created. The widget controls a physical Raman microscopy rig; verify stage travel,
focus range, shutter state, and sample clearance before running any automated motion.
\end{abstract}
\vfill
{\footnotesize Documentation edition based on the public source; editorial revision by Jason W. Yu.}
\clearpage
\tableofcontents

\chapter{Architecture and Workflow Dependencies}

\section{Main components}
The application combines napari visualization, Qt controls, Micro-Manager device
control, Raman acquisition, coordinate transformation, autofocus, calibration,
cell selection, stage positioning, and multidimensional acquisition. The sidebar
organizes these capabilities into a sequence of expandable sections. Plot windows
show spectra and calibration results, while the log window records acquisition progress.

Begin each hardware session with \ui{Connect hardware}. The bottom status label
reports the current action and confirms when each workflow is ready.

\section{The selection-to-MDA contract}
\ui{Run Raman MDA} uses three prepared variables: \code{sources},
\code{autofocus\_p}, and \code{new\_seq}. Prepare these variables with one of the
following selection workflows:
\begin{itemize}
  \item \ui{Run automated selection};
  \item \ui{Manual selection} (after the user clicks points in the created layers);
  \item \ui{Center clicked cells} (optional conversion of a non-batch manual selection);
  \item \ui{Generate grid} in the stage-grid section.
\end{itemize}

Before selection begins, live acquisition stops and the napari-micromanager MDA
settings are prepared with axis order \code{t,p,c,z} and a centered Z plan. The
recommended sequence is:
\begin{enumerate}
  \item Connect hardware; this attempts to open the \code{napari-micromanager} dock.
  \item In that MDA widget, define or verify stage positions and the basic sequence.
  \item Use automated selection, manual selection, or stage-grid generation to
        create the source layers and revised sequence.
  \item For manual selection, finish clicking the intended cells in the layers.
  \item Configure \ui{Run Raman MDA}, then start it.
\end{enumerate}
The MDA dock provides the underlying position and sequence information used by
the selection controls.

\section{Coordinate conventions}
Point and shape workflows use the \emph{last} napari layer. Point workflows use
the first point in that layer and its final two coordinates. For a point stored as
(y, x), the transformer converts the normalized bright-field coordinate into
galvanometer volts. Place the intended Points or Shapes layer last in the layer list.

\chapter{Loading}
This is the only section expanded at startup. Choose paths and the output folder
before connecting, because connection loads the devices and changes the process
working directory.

\fieldtable
Micro-Manager config (.cfg) & \code{test3.cfg} & Device configuration used for
the device setup. The browse button filters for CFG files. During
connection the widget also attempts a GFP-to-BF channel warm-up.\\
Transformer model (.json) & blank; example \code{model\_2026-01-08.json} &
Coordinate transformer used for bright-field pixel to Raman galvo conversion.
Point collection, calibration, reference acquisition, mapping, and Raman MDA require it.\\
Vandermonde model (.json) & blank; example \code{vandermonde\_model.json} &
Pixel-to-stage model used only when cells are physically centered. Required by
automated \ui{Center cell} and \ui{Center clicked cells}. It complements the Raman
aiming transformer.\\
Output folder & blank = current directory & Created if needed, then made the process
working directory during connection. Relative reference, map, model, and MDA paths
are resolved from here. Changing the text after connection does not change directory.\\
Center wavelength (nm) & 785.00; 0--2000 & Desired spectrometer center wavelength.
The bold label shows the hardware value. \ui{Update wavelength} is enabled only after
connection and applies the displayed wavelength to the spectrometer.\\
Grating & populated after connection & Lists the spectrometer gratings using 1-based
indices. \ui{Update grating} moves the turret and reports groove density and center
wavelength. Wait for motion to finish before acquiring.\\
\endfieldtable

\ui{Connect hardware} unloads old devices, loads the CFG, opens the
napari-micromanager dock, creates the spectrum collector/DAQ, loads both models,
and refreshes channels, wavelength, and gratings. \ui{Reload transformer} reloads
the Raman transformer and Vandermonde model without reconnecting. \ui{Disconnect}
unloads devices and clears all calibration, selection, writer, and model state;
selection is prepared again after reconnecting.

\section*{Saved paths}
The CFG, transformer, and Vandermonde fields point to existing input files. The
\ui{Output folder} becomes the base location for relative results. For example,
\code{data/run} is resolved as \code{<Output folder>/data/run}.

\section*{Options at a glance}
\begin{longtable}{@{}p{0.31\textwidth}p{0.63\textwidth}@{}}
\toprule\textbf{Option} & \textbf{Brief definition}\\\midrule
Micro-Manager config & Hardware and device configuration file.\\
Transformer model & Bright-field-pixel to Raman-galvo coordinate model.\\
Vandermonde model & Image-pixel to stage-motion model for centering cells.\\
Output folder & Working directory applied when hardware connects.\\
Center wavelength & Spectrometer center wavelength setting.\\
Grating & Installed spectrometer grating selection.\\
Connect / Disconnect & Initializes or unloads the rig and session state.\\
Reload transformer & Reloads the coordinate models without reconnecting.\\
\bottomrule\end{longtable}

\chapter{Collect Spectra Using Points Layer}

\textbf{Preparation:} connect the collector and DAQ, load the transformer, and
place a Points layer last. The first stored point is used for acquisition. Keeping
one point in this layer makes the intended target clear.

\fieldtable
Exposure (ms) & 1000; 1--1,000,000 & Integration time used by the spectrum collector.\\
N (repeats, at least 2) & 2; 2--1000 & Number of identical copies of the transformed
galvo coordinate used for acquisition. The minimum of two is enforced.\\
Save as & blank = display only & Optional filename. If the \code{.npy} suffix is
provided, the widget adds it. A relative name is saved under the active working directory.\\
\endfieldtable

\ui{Collect spectra at last point} restarts the galvo, transforms the point,
collects N spectra, optionally saves the NumPy array, and opens a spectrum plot.

\section*{Saved output}
When \ui{Save as} contains a name, spectra are saved as
\code{<output folder>/<name>.npy}. With a blank name, the spectra are displayed
without creating a data file.

\section*{Options at a glance}
\begin{longtable}{@{}p{0.31\textwidth}p{0.63\textwidth}@{}}
\toprule\textbf{Option} & \textbf{Brief definition}\\\midrule
Exposure & Integration time for each Raman acquisition.\\
N (repeats) & Spectra acquired at the same transformed point.\\
Save as & Optional NumPy filename; blank disables saving.\\
Collect spectra & Transforms the point, acquires spectra, and displays them.\\
\bottomrule\end{longtable}

\chapter{Laser Aiming Calibration and Recalibration}

\fieldtable
N (repeats) & 20; 1--1000 & Spectra acquired per calibration target. More repeats
increase acquisition time and can stabilize detection.\\
Exposure (ms) & 1000; 1--1,000,000 & Raman exposure used by the calibrator.\\
Max volts & 1.80; 0.01--10.00 & Galvanometer-voltage extent used by the
calibration. Set it within the confirmed operating range of the rig.\\
Grid size & 10; 2--100 & Controls calibration
sampling density. Larger grids require substantially more measurements.\\
Threshold & 1.00; 0--100 & Detection threshold used to accept or localize
calibration responses.\\
Recalibration & off & Reveals the controls for correcting the current calibration.\\
Model name & blank & Base filename for the corrected transformer. The save routine
writes \code{<name>.json}, loads it immediately, and updates the Loading path.\\
\endfieldtable

\ui{Run calibration} uses the active transformer and stores the resulting dataset
in memory. It opens a log and calibration plot. Recalibration uses the most recent
calibration dataset produced in the connected session.

For manual recalibration, check \ui{Recalibration}, click \ui{Open manual selector},
and then work through the displayed calibration frames: click a point; press Enter
to advance; Backspace to return; R to reset; or N to mark a frame as NaN. Close or
finish the selector, enter \ui{Model name}, and click \ui{Save recalibrated model}.

\section*{Saved output}
The initial calibration remains in session memory for plotting and recalibration.
\ui{Save recalibrated model} writes
\code{<output folder>/<Model name>.json} and immediately activates that model.

\section*{Options at a glance}
\begin{longtable}{@{}p{0.31\textwidth}p{0.63\textwidth}@{}}
\toprule\textbf{Option} & \textbf{Brief definition}\\\midrule
N (repeats) & Spectra used at each calibration target.\\
Exposure & Raman integration time during calibration.\\
Max volts & Galvanometer-voltage extent of the calibration.\\
Grid size & Sampling density of the calibration grid.\\
Threshold & Response criterion used by the calibrator.\\
Run calibration & Acquires and plots a new calibration dataset.\\
Recalibration & Reveals the manual correction controls.\\
Model name & Base filename for the corrected transformer JSON.\\
Open manual selector & Opens the last calibration dataset for correction.\\
Save recalibrated model & Saves, loads, and activates the corrected model.\\
\bottomrule\end{longtable}

\chapter{Axial Background Scan}

This section collects an autofocus/background Z series at one Raman point. Put a
single-point Points layer last; the first point is transformed to galvo volts and
repeated N times at every Z sample.

\fieldtable
Name & required; example \code{testing} & Human-readable prefix for the output.\\
Exposure (ms) & 1000; 1--1,000,000 & Exposure for each Raman spectrum.\\
N (spectra per z) & 5; 1--1000 & Replicate spectra at each axial position.\\
Search range (um) & 10; 0.1--1000 & Half-range around the starting Z; displayed
coordinates span negative range to positive range. Verify objective clearance.\\
Search pts & 20; 2--500 & Number of axial samples across the full search interval.\\
\endfieldtable

\ui{Collect reference spectra} performs the autofocus/background scan, moves the
stage to the selected focus Z, plots all spectra, and writes
\code{reference/<name>\_<uuid>.zarr}. The dataset contains \code{spec[z,n,pixel]}
plus Z, point X/Y, and exposure metadata. The stage is intentionally left at the
found focus position.

\section*{Saved output}
Each run is saved as
\code{<output folder>/reference/<Name>\_<uuid>.zarr}. The Zarr dataset stores the
spectra, Z coordinates, selected point coordinates, and exposure value.

\section*{Options at a glance}
\begin{longtable}{@{}p{0.31\textwidth}p{0.63\textwidth}@{}}
\toprule\textbf{Option} & \textbf{Brief definition}\\\midrule
Name & Prefix identifying the saved reference dataset.\\
Exposure & Integration time for each reference spectrum.\\
N (spectra per Z) & Replicate spectra acquired at every Z position.\\
Search range & Positive and negative axial search extent.\\
Search points & Number of evenly spaced Z samples.\\
Collect reference spectra & Runs autofocus/background collection, plots, and saves.\\
\bottomrule\end{longtable}

\chapter{Spatial Mapping}

Create a Shapes layer, draw a rectangle, and make that layer last. The code uses
only the first shape and reduces it to its min/max bounds; a rotated rectangle or
nonrectangular polygon is therefore scanned as its axis-aligned bounding box.

\fieldtable
File name & required; example \code{scan\_label} & Label inserted in the Zarr filename.\\
Raman exposure (ms) & 1000; 1--1,000,000 & Exposure at every Raman grid coordinate.\\
N (grid side) & 20; 2--500 & Produces an N by N grid, hence N squared Raman points
per Z plane. Doubling N approximately quadruples point count.\\
Z offset (um) & 4; -1000 to 1000 & The base Raman Z is current Z minus this
offset. The stage returns to its original Z after a successful scan.\\
Z-scan & off & When enabled, collects the full Raman grid at multiple Z planes.\\
Z range (+/- um) & 5; 0.1--1000 & Half-range around the base Raman Z; visible only
for Z-scan.\\
Z steps & 10; 2--500 & Evenly spaced planes across the full Z range.\\
Extra channel & none; 1--1,000,000 ms & \ui{+ Add channel} adds a Micro-Manager
channel and exposure row; \ui{x} removes it. Duplicate channels are ignored. BF is
excluded here because BF is always acquired before and after the scan.\\
\endfieldtable

At start the widget snaps BF at 10 ms and one image for each extra channel. It then
switches to RM, opens \code{Fluoshutter}, collects the Raman grid, closes the shutter,
restores Z, and snaps ending BF. A Z-scan also records BF at every plane. Outputs are
\code{grid\_scan\_data\_<label>\_<uuid>.zarr} or
\code{grid\_scan\_z\_<label>\_<uuid>.zarr}, followed by a plot window.

\section*{Saved output}
A single-plane map is saved as
\code{<output folder>/grid\_scan\_data\_<File name>\_<uuid>.zarr}. A multi-plane
map is saved as
\code{<output folder>/grid\_scan\_z\_<File name>\_<uuid>.zarr}.

\section*{Options at a glance}
\begin{longtable}{@{}p{0.31\textwidth}p{0.63\textwidth}@{}}
\toprule\textbf{Option} & \textbf{Brief definition}\\\midrule
File name & Label inserted into the grid-scan filename.\\
Raman exposure & Integration time at every grid coordinate.\\
N (grid side) & X and Y sample count; total is N squared per plane.\\
Z offset & Raman-plane displacement from the current Z.\\
Z-scan & Enables multi-plane Raman acquisition.\\
Z range / Z steps & Axial half-range and number of sampled planes.\\
Add/remove channel & Manages fluorescence-channel and exposure pairs.\\
Run grid scan & Acquires BF, fluorescence, Raman grid data, and ending BF.\\
\bottomrule\end{longtable}

\chapter{Generate Stage Grid}

This section creates stage positions around the \emph{current} stage XY and stores
the prepared variables for \ui{Run Raman MDA}.
Every position carries the same fixed image point.

\fieldtable
Autofocus object & None, laser, software, quartz, glass, cell & Autofocus mode attached
to the generated positions. The choice also controls which MDA autofocus fields are visible.\\
FOV x / FOV y (px) & 740 / 540; 0--100000 & Fixed image coordinate used at every
field of view. Note the UI labels X then Y; internal selection receives both separately.\\
X range / Y range (+/- um) & 100 / 100 & Half-width of the stage grid about the
current XY. Total spans are twice these values.\\
X step / Y step (um) & 50 / 50; minimum 0.01 & Stage-position spacing. Confirm
that all generated positions remain within stage and sample limits.\\
Repeats (points per position) & 2; 2--1000 & Identical points at each stage position;
minimum two is required by the DAQ.\\
Skip BF pre-scan & on & Uses blank placeholder images to establish dimensions. Turn
off when real per-position BF images are needed for setup.\\
\endfieldtable

\ui{Generate grid} stops live mode, prepares the napari-micromanager MDA sequence,
builds the stage grid, and prepares \code{sources}, \code{autofocus\_p}, and
\code{new\_seq}. After the status reports ``stage grid ready,'' proceed to
\ui{Run Raman MDA}.

\section*{Saved output}
This section prepares variables in memory. The stage positions and Raman results
are written later by \ui{Run Raman MDA} to its \ui{Writer output dir}.

\section*{Options at a glance}
\begin{longtable}{@{}p{0.31\textwidth}p{0.63\textwidth}@{}}
\toprule\textbf{Option} & \textbf{Brief definition}\\\midrule
Autofocus object & Focus strategy attached to generated positions.\\
FOV X / FOV Y & Fixed pixel target repeated at every position.\\
X/Y range & Grid half-widths around the current stage XY.\\
X/Y step & Stage spacing between neighboring positions.\\
Repeats & Identical aiming points carried by every position.\\
Skip BF pre-scan & Uses blank placeholders instead of BF setup images.\\
Generate grid & Creates positions and selection results for Raman MDA.\\
\bottomrule\end{longtable}

\chapter{Automated Cell Selection}

\section{Mask controls}
\fieldtable
Center Y / Center X & 540 / 740; 0--100000 & Center of the circular permitted region
in image pixels, in Y,X order.\\
Radius & 100; 1--100000 & Radius of the permitted circular region. The same values
are passed later to the Raman MDA engine.\\
Add mask & button & Adds a red masked overlay, with a green 10-pixel center marker,
to help inspect the permitted region before selection.\\
Click to center & toggle button & Arms a one-shot viewer click that moves the stage
so the clicked feature reaches Center Y/X. Requires a loaded Vandermonde model and
connected hardware; use cautiously because it commands stage motion.\\
\endfieldtable

\section{Selection controls}
\fieldtable
Autofocus object & laser; choices laser, software, quartz, glass, cell, None &
Autofocus strategy saved with the selection. None disables autofocus in the later MDA.\\
N per FOV & 6; 1--1000 & Requested cell count. Automated selection passes
\code{N+1} to its helper. In manual batch mode it is the exact click count per FOV.\\
Center cell & off & Splits detections into one new stage position per cell and shifts
each cell to the requested center. Requires the Loading Vandermonde model.\\
Aiming pattern & Square or Circle & Shape of the Raman subpoint pattern placed around
each selected cell.\\
Pattern size (px) & 30; 0--10000 & Spatial extent of the point pattern.\\
N\_x (subpoints) & 1; 1--100 & Pattern sampling parameter. Batch MDA requires the
resulting pattern multiplier to be at least two.\\
Background distance (px) & 80; 0--1,000,000 & Distance threshold used by automated selection.\\
Integrated batch collection & False & In batch mode positions/click counts and Raman
point handling differ; it must agree with dataset generation.\\
Cellpose model & installed models; defaults to cyto2 if available & Segmentation model
used to identify candidate cells.\\
\endfieldtable

\ui{Run automated selection} prepares the external MDA widget, performs Cellpose-based
selection, creates source layers and a new sequence, and stores them for Raman MDA.

\ui{Manual selection} creates empty point-source layers. In batch mode, click exactly
\ui{N per FOV} cells per FOV; in non-batch mode, click freely. The button immediately
stores references to those layers, so finish clicking before starting the Raman MDA.
For non-batch manual results, \ui{Center clicked cells} uses the Vandermonde model to
turn each clicked cell into a centered stage position, replacing the selection results.

\section*{Saved output}
Automated and manual selection prepare \code{sources}, \code{autofocus\_p}, and
\code{new\_seq} in memory. \ui{Add mask} adds a napari layer in the current viewer.
The subsequent Raman MDA writes the acquired data to \ui{Writer output dir}.

\section*{Options at a glance}
\begin{longtable}{@{}p{0.31\textwidth}p{0.63\textwidth}@{}}
\toprule\textbf{Option} & \textbf{Brief definition}\\\midrule
Center Y / Center X & Pixel center of the permitted mask and centering target.\\
Radius & Radius of the circular permitted region.\\
Add mask & Displays the configured circular selection region.\\
Click to center & Moves a clicked feature to the configured center.\\
Autofocus object & Focus strategy saved with the selected positions.\\
N per FOV & Requested detections or manual batch clicks per field.\\
Center cell & Produces one centered stage position per detected cell.\\
Aiming pattern / size / N\_x & Shape, extent, and sampling of Raman subpoints.\\
Background distance & Distance threshold for automated selection.\\
Integrated batch collection & Selects batch or individual point handling.\\
Cellpose model & Segmentation model used to detect cells.\\
Run automated selection & Segments cells and prepares MDA selection results.\\
Manual selection & Creates layers for hand-clicked cell selection.\\
Center clicked cells & Converts non-batch clicks into centered positions.\\
\bottomrule\end{longtable}

\chapter{Run Raman MDA}

\textbf{Preparation:} connect with both collector and transformer, then run
automated selection, manual selection, centered-manual selection, or stage-grid
generation. These workflows prepare the variables used by the MDA.

\fieldtable
Writer output dir & \code{data/run} & Directory where Raman TIFF and NumPy outputs
are written; blank uses \code{data/run}.\\
Autofocus positions & blank = selection & Optional comma-separated zero-based stage
position indices, e.g. \code{0,2,5}. \code{None} behaves like blank. Out-of-range
indices are rejected.\\
Imaging positions & blank = autofocus positions & Optional comma-separated zero-based
indices at which imaging occurs. The code initializes this from the original
selection autofocus list before applying overrides, so set this explicitly when
using a substantially different autofocus override.\\
Raman glass offset (um) & 5; -1000 to 1000 & Axial offset used for the Raman measurement.\\
Autofocus search range & 6; 0.1--1000 & Coarse autofocus range; hidden if autofocus is None.\\
Autofocus search pts & 8; 2--500 & Coarse autofocus sample count.\\
Laser fine search range (+/- um) & 1.5; 0.1--1000 & Fine range shown only for laser autofocus.\\
Laser fine search pts & 8; 2--500 & Fine laser-autofocus sample count.\\
Segment and track & off & Re-segments images and updates aiming during a time series.\\
Segment channel & BF initially & Micro-Manager channel used for segmentation.\\
Image rescale factor & 2.0; 0.1--100 & Image scale used during segmentation.\\
Cellpose model & installed models; cyto2 preferred & Model for time-series re-segmentation.\\
Crop image around mask & True & Whether segmentation is cropped to the circular mask region.\\
Tracking config (.json) & \code{particle\_config.json} & Particle-tracking configuration.\\
Exposure per cell (ms) & 1000; 1--1,000,000 & Total exposure supplied while building
the Raman sequence. In non-batch mode it is multiplied by the aiming-pattern multiplier.\\
Loops (time points) & 100; 1--1,000,000 & Number of temporal repetitions.\\
Interval (s) & 600; 0--1,000,000 & Requested interval between time points.\\
Refocus \& segment every & 1 & Cadence in time points; 1 means every time point.\\
Z relative (comma-separated um) & \code{0, 4} & Relative Z planes used to replace
the sequence Z plan. Values must parse as floats.\\
Raman z indices & \code{0} & Zero-based indices into the Z list where Raman is
requested. Ensure every index exists; e.g. with two Z values, valid indices are 0 and 1.\\
Extra fluorescence channel & none; default row exposure 10 ms & \ui{+ Add channel}
adds any Micro-Manager channel (including BF); \ui{x} removes it. Duplicate channel
names are ignored. Added channels inherit the first sequence channel as a template.\\
\endfieldtable

The acquisition uses axis order \code{t,p,c,z}, assigns the selection's aiming-source
layers, replaces the time and Z plans, adds unique fluorescence channels, and writes
Raman metadata containing the requested Z indices. \ui{Stop} requests MDA cancellation
and stops sequence acquisition; the current hardware event may finish before exit.

\section{Generate dataset}
\ui{Generate dataset} asks for a completed MDA run directory. It loads the experiment,
using the current \ui{Integrated batch collection} value, and writes sibling outputs
\code{data/dataset/ds\_<run>.zarr} and \code{df\_<run>.pkl}. It also opens the dataset
viewer. Set the batch dropdown to match how that run was acquired before generating.

\section{Pixel-to-stage calibration}
Check \ui{Pixel-to-stage calibration} to reveal these controls:
\fieldtable
Dataset (.zarr) & required & Generated dataset containing a JSON
\code{useq\_sequence} attribute and one image per stage position.\\
Vandermonde degree & 1; 1--5 & Polynomial degree. Degree d needs at least
(d+1)(d+2)/2 valid picked points. Prefer the lowest degree with adequate residuals.\\
Pick points & button & Opens images position by position. Click the same recognizable
feature in each frame; skipped/NaN frames are excluded.\\
Model file & \code{vandermonde\_model.json} & Suggested save name. A save dialog still
asks for the final location.\\
Fit \& save model & button & Centers image and stage coordinates, reports degree 1--3
RMSE where possible, fits the selected degree, and saves JSON. The saved path is also
copied into Loading's Vandermonde field for immediate center-cell use.\\
\endfieldtable

\section*{Saved output}
The Raman acquisition is saved under
\code{<output folder>/<Writer output dir>/}. The writer creates the TIFF and NumPy
files for the run. \ui{Generate dataset} saves
\code{<run parent>/dataset/ds\_<run>.zarr} and
\code{<run parent>/dataset/df\_<run>.pkl}. Pixel-to-stage calibration saves the
chosen \ui{Model file} as JSON at the location selected in the save dialog.

\section*{Options at a glance}
\begin{longtable}{@{}p{0.31\textwidth}p{0.63\textwidth}@{}}
\toprule\textbf{Option} & \textbf{Brief definition}\\\midrule
Writer output directory & Destination used by the Raman TIFF/NumPy writer.\\
Autofocus positions & Zero-based position indices that receive autofocus.\\
Imaging positions & Zero-based position indices that receive imaging.\\
Raman glass offset & Axial offset between glass focus and Raman measurement.\\
Autofocus range / points & Extent and sampling count of the coarse focus search.\\
Laser fine range / points & Extent and sampling count of laser fine focus.\\
Segment and track & Re-detects cells and updates aiming during the time series.\\
Segment channel & Image channel used for segmentation.\\
Image rescale factor & Resampling factor supplied to segmentation.\\
Cellpose model & Model used for repeated cell segmentation.\\
Crop around mask & Restricts segmentation to the configured mask region.\\
Tracking config & JSON settings used for cell/particle tracking.\\
Exposure per cell & Raman exposure budget assigned to each selected cell.\\
Loops / interval & Number of time points and requested spacing between them.\\
Refocus \& segment every & Time-point cadence for focus and tracking updates.\\
Z relative & Comma-separated relative Z planes in micrometres.\\
Raman Z indices & Indices of relative Z planes receiving Raman acquisition.\\
Add/remove fluorescence channel & Manages channel/exposure pairs in the sequence.\\
Run Raman MDA & Builds and launches the final \code{t,p,c,z} acquisition.\\
Stop & Requests MDA cancellation and stops sequence acquisition.\\
Generate dataset & Converts a completed run into Zarr and pickled-table outputs.\\
Pixel-to-stage calibration & Reveals the Vandermonde calibration workflow.\\
Dataset & Zarr images and stage coordinates used for calibration.\\
Vandermonde degree & Polynomial degree used for the pixel-to-stage fit.\\
Pick points & Opens the common-feature picker for each stage position.\\
Model file & Suggested JSON filename for the fitted model.\\
Fit \& save model & Fits, reports residuals, saves, and activates the model.\\
\bottomrule\end{longtable}

\chapter{Session Checklist and Output Summary}
\begin{enumerate}
  \item Verify the CFG, Raman transformer, optional Vandermonde model, and output folder.
  \item Connect; confirm the status, wavelength, grating, and available channels.
  \item Verify the intended layer is last before point or shape operations.
  \item For MDA work, configure the napari-micromanager MDA widget and run exactly one
        selection-producing workflow. Finish manual clicks before continuing.
  \item Check all physical motion ranges, point counts, Z planes, loops, and exposures.
  \item Run a small pilot, inspect the log and plots, then scale to the full acquisition.
\end{enumerate}

\begin{longtable}{@{}p{0.32\textwidth}p{0.62\textwidth}@{}}
\toprule\textbf{Workflow} & \textbf{Output}\\\midrule
Point spectrum & Optional \code{<name>.npy}; plot window.\\
Calibration & In-memory dataset and plot; recalibration saves \code{<model>.json}.\\
Axial background scan & \code{reference/<name>\_<uuid>.zarr}.\\
Spatial mapping & \code{grid\_scan\_data\_*.zarr} or \code{grid\_scan\_z\_*.zarr}.\\
Selections / stage grid & In-memory source layers, autofocus indices, and MDA sequence.\\
Raman MDA & TIFF/NumPy writer directory under \ui{Writer output dir}.\\
Generate dataset & \code{ds\_<run>.zarr} and \code{df\_<run>.pkl}.\\
Pixel-to-stage fit & Vandermonde JSON model.\\\bottomrule
\end{longtable}

Source: \url{https://github.com/JasonYu1/napari-raman-widget}. The project is
distributed under the BSD 3-Clause license.
\end{document}
